\subsection{File System Structure and Organization} \label{file-structure-and-hierarchy} % done

FreeBSD natively supports two filesystems: Unix File System 2 (UFS2) and Zettabyte File System (ZFS).
The \textbf{Unix File System (UFS)} is a family of filesystems designed for Unix systems. UFS2 extends the original UFS
with some major improvements, which includes the addition of 64-bit block pointers, enhancements for better support of
variable-size blocks, flag fields extensions, per-inode extended attribute extent additions, and the implementation
of lazy inode initialization \parencite{van_gelderen2003}. It was adopted in FreeBSD 5.0 and has been refined for around
two decades ever since \parencite{freebsd_zfs_vs_ufs}. UFS2 partitions each disk slice into \textit{cylinder groups},
which are composed of \textit{disk cylinders} \parencite{oracle2010}. Each disk cylinder is further subdivided into
\textit{logical blocks} assigned with an address to manage file organization within the cylinder group \parencite{oracle2010}.
There are four types of blocks according to \textcite{oracle2010}:

\begin{itemize}
	\item \textbf{Boot Block}: The boot block stores the data necessary in booting the operating system.
	      The boot block is typically located at the first cylinder group or \textit{group 0} and occupies around 8K in
	      a data slice. When the filesystem is not used for booting the system, the boot block is left blank

	\item \textbf{Superblock}: The superblock stores information about the filesystem, which can include the size and
	      status of the filesystem, the name of the filesystem and its disk volume, the size of the logical blocks, when the
	      filesystem is last updated, the size of the cylinder groups, how many data blocks are stored in a cylinder group,
	      the current state of the filesystem, and the path to the last mounted device. The superblock is typically located
	      at the start of each disk slice. The superblock is often replicated in a cylinder group since it contains critical
	      information crucial in managing the filesystem. Whenever the superblock or the disk slice holding the superblock
	      gets corrupted or damaged, the information about the filesystem can still be retrieved from its replicas.
	      The replicas are typically offsetted by a specific amount depending on the type of the disk. For multi-platter disk
	      drives for instance, the offsets are configured such that a replica of the superblock is present in each platter.
	      The superblock also contains a summary information block, which notes the changes in the filesystem and the number
	      of inodes, directories, and storage blocks that are currently used in the filesystem. The summary information
	      block is not replicated, hence it can only be found with the main superblock.

	\item \textbf{Inode}: The inode stores information about a file, which can include the type of the file, the
	      permissions set on the file, hard links to the file, the owner of the file, the group where the file belongs,
	      the file size, disk-block addresses where the file is stored, and when the file was last created, modified,
	      and accessed. It is important to note that the inode does not store the name of the file, as it is commonly
	      stored in the directory the file is at.

	\item \textbf{Storage/Data Block}: The storage block contains the data associated with a file. The sizes of the
	      storage blocks are known when the filesystem is first created. Storage blocks can be allocated to either a block
	      size of 8K or a fragment size of 1K. Storage blocks can either store the contents of a file of the inode
	      entries with their corresponding file names of a directory.
\end{itemize}

UFS2 is commonly used in older versions of FreeBSD. In newer versions, however, the \textbf{Zettabyte File System (ZFS)}
is widely used. One key characteristic of ZFS is that it is both a filesystem and a volume manager \parencite{freebsd_zfs},
which allows it to organize files and manage disks in one system.
The ZFS filesystem is composed of \textit{Data Management Unit (DMU)} objects \parencite{siebenmann2018},
which can represent a file, directory, disk volume, or internal metadata \parencite{openzfs_docs_dnode}. Each DMU object
is defined by a \textit{dnode}, which holds the metadata of the object \parencite{openzfs_docs_dnode}. DMU objects are represented
as a tree of block pointers on disk \parencite{openzfs_docs_dnode}. The root node is a \texttt{dnode\_phys\_t} structure
which contains the object metadata. The metadata typically consists of the object type, the number of levels in the
tree hierarchy of the object, and the amount of memory space used by the object \parencite{openzfs_docs_dnode}.
ZFS also utilizes a \textit{Dataset and Snapshot Layer (DSL)} which manages the filesystems and its clones and snapshots
\parencite{openzfs_docs_dnode}. The DSL typically consists of DSL directories and datasets which are stored as dnodes
in a special object set called the \textit{Meta Object Set} \parencite{openzfs_docs_dnode}. Each DSL directory or dataset
in a tree in itself \parencite{openzfs_docs_dnode}.

In NetBSD, the Fast File System (FFS), Log-Structured File System (LFS), and Chip File System (CHFS) were officially adopted
on different versions of NetBSD.
Both \textbf{Fast File System (FFS)} and \textbf{Log-Structured File System (LFS)} share a lot of similarities with UFS in
terms of its organization. For instance, the filesystem subdivides the disk into \textit{cylinder groups}
\parencite{arpaci_dusseau2023}. Each cylinder in a cylinder group is a set of tracks that share the same distance
from the center \parencite{arpaci_dusseau2023}. Similar to UFS, each cylinder group contains a superblock with its
replicas, a list of inodes, and the data blocks \parencite{arpaci_dusseau2023}. Additionally, the filesystem stores
the inode and data bitmaps before the inodes and data blocks. The bitmaps essentially keep track of where the inodes
and data blocks are allocated in the group. In LFS, each group or segment contains a summary block just like with UFS
\parencite{rosenblum1991}. However, summary blocks can be chained together in a linked list, with a summary block
holding a pointer to the next summary block \parencite{rosenblum1991}. The filesystem then read these summary block
as one long linear log \parencite{rosenblum1991}. NetBSD also adopts the \textbf{Chip File System (CHFS)},
which was primarily designed for flash devices \parencite{chfs_about2018}. CHFS utilizes an intermediary layer between
the filesystem and the media device called the \textit{eraseblock-handler (EBH)}. EBH essentially maps the physical eraseblocks
(PEB) of the device to the logical eraseblocks (LEB) managed by the filesystem.

It is evident that the approaches of FreeBSD's ZFS and NetBSD's FFS in file management are largely different. NetBSD's FFS
follows the standard approach of using cylinder groups in a single pre-partitioned disk slice. FreeBSD's approach is rather
different by unifying file and disk management into a single system. An advantage of FFS is that it exhibits high cross-platform
compatibility primarily because of its simplicity. The \textit{dnode/object set/DSL} hierarchy model of ZFS is complex.
Considering the additional features that ZFS comes with, it might not be well-supported on constrained platforms or
architectures.

\subsection{File Access Methods and Operations} \label{file-access-methods} % done

In FreeBSD, the VFS API provides two kinds of operations for working with files: \textit{path-based operations} and
\textit{file-descriptor operations} \parencite{jensen2021}. The \texttt{open()} function can be used to retrieve a file
descriptor given a path when the \texttt{O\_PATH} flag is specified \parencite{jensen2021}.
\texttt{open()} returns a non-negative integer if the operation is
successful \parencite{freebsd_open}. The returned integer corresponds to the file descriptor, which is the lowest-numbered
file descriptor that is currently not being used by the system \parencite{freebsd_open}. On failure, it returns \texttt{-1}
\parencite{freebsd_open}. File descriptors are managed by the \texttt{filedesc} struct \parencite{freebsd_filedesc_h}.
The struct contains \texttt{fd\_files} of type \texttt{struct fdescenttbl*}, which is a table of files that are currently
opened; \texttt{fd\_freefile} of type \texttt{int}, which points to the next approximately free file in the files table; and
\texttt{fd\_files->fd\_nfiles}, which tracks the current number of allocated slots in the files table \parencite{freebsd_filedesc_h}.
File descriptors can be referred by a process with the \texttt{proc} struct's \texttt{p\_fd} attribute when a process is opening a
file \parencite{freebsd_proc_h}. FreeBSD supports a variety of file descriptor data types. File descriptors of files have a
type of \texttt{DTYPE\_VNODE} \parencite{freebsd_file_h}. Other file descriptor types include \texttt{DTYPE\_SOCKET} for sockets,
\texttt{DTYPE\_PIPE} for pipes, \texttt{DTYPE\_MQUEUE} for POSIX message queues, \texttt{DTYPE\_SEM} for POSIX semaphores,
\texttt{DTYPE\_DEV} for device-specific file descriptors, and \texttt{DTYPE\_PROCDESC} for process descriptors \parencite{freebsd_file_h}.
File descriptors are typically used using the \texttt{AT\_EMPTY\_PATH} flag, which is specified on the \texttt{at()} set of syscalls
\parencite{jensen2021}. The syscalls can then perform operations on the file specified by its file descriptor when passed on the
\texttt{dirfd} argument and the path argument is set to an empty string literal \parencite{jensen2021}.

NetBSD supports high-level operations for working with vnodes. For instance, the function \texttt{vn\_open()} is used for common
open operations on the vnode \parencite{netbsd_vnopen}. For vnode close operations, the function \texttt{vn\_close()} is used
\parencite{netbsd_vnopen}. \texttt{vn\_open()} can either trigger the \texttt{VOP\_CREATE()} operation, which creates a new file,
or the \texttt{VOP\_OPEN()}, which opens an existing file, depending on the flags passed \parencites{netbsd_vnopen}{netbsd_vopopen}.
\texttt{vn\_open()} can either return a vnode or a file descriptor number \parencite{netbsd_vnopen}. If it returns a vnode, it locks the
vnode and is returned in the \texttt{ret\_vp} argument \parencite{netbsd_vnopen}. If it returns a file descriptor number, it is returned
instead in the \texttt{ret\_fd} argument \parencite{netbsd_vnopen}. For file access, the \texttt{open()} and \texttt{openat()}
functions are implemented for opening or creating files, passing in an absolute path for \texttt{open()} and a relative path for
\texttt{openat()} \parencite{netbsd_open}. Files can be opened in either read mode, write mode, execute mode, or reading and writing
mode with the \texttt{O\_RDONLY}, \texttt{O\_WRONLY}, \texttt{O\_EXEC}, and \texttt{O\_RDWR} respectively \parencite{netbsd_open}.
Additional flags can be passed, which include \texttt{O\_APPEND} for appending a file, \texttt{O\_CREAT} for creating a file if
it does not exist yet, and \texttt{O\_DIRECTORY} for accessing directories \parencite{netbsd_open}. All files that are currently
being used are being referenced in a file descriptor table \parencite{netbsd_file}. Each entry in the file descriptor table
holds a count of references with the \texttt{f\_count} attribute in the \texttt{file} struct \parencite{netbsd_file}. The count gets
incremented when new references are added and decremented when the \texttt{close()} function is called \parencite{netbsd_file}.
The file is eventually freed when the references count reaches 0 \parencite{netbsd_file}. The \texttt{f\_offset} attribute specifies
the number of bytes added for each read or write operation \parencite{netbsd_file}. This offset cannot be stored in the file's inode
since multiple processes can open the same file each with their own offsets \parencite{netbsd_file}. Write operations are handled
by the \texttt{vn\_write()} function \parencite{netbsd_vnwrite}. When the flag \texttt{FOF\_UPDATE\_OFFSET} is set, the position
will be updated when read or write operations are completed \parencite{netbsd_vnwrite}. Files can be mapped into their respective
sectors in memory with the \textit{Unified Buffer Cache (UBC)} \parencite{netbsd_ubc}. The mapping operation is done using \texttt{mmap()}
\parencite{netbsd_mmap}. Whenever a mapping operation will be performed, the \texttt{VOP\_MMAP} function is called to notify the
filesystem \parencite{netbsd_vnodeops}.

Both systems adopt the vnode abstraction as the primary representation of an open file to the kernel. Both also adopts the
POSIX-compliant operations for accessing open files such as \texttt{open()}, \texttt{read()}, \texttt{write()}, and \texttt{close()}.
Both FreeBSD and NetBSD also share the same mechanisms in handling file accesses from multiple processes. Both of them also takes
a modular approach in extending the static and dynamic sections of vnodes. However, FreeBSD has a richer
model for representing file descriptors, which can support multiple file descriptor types. This allows multiple resources to be
handled on a unified descriptor table. This is unlike NetBSD that can only access and manage filesystem objects represented as
vnodes.

\subsection{Directory Structures} \label{dir-struct} % done

The filesystems in both FreeBSD and NetBSd adopts a hierarchical tree-based structure reminiscent to UNIX filesystems. The root directory of the
filesystem is the \texttt{/} directory \parencite{freebsd_dir_struct}. Table \ref{root-subdirs} lists the common subdirectories of the
root directory according to \textcite{freebsd_dir_struct}. Additional filesystems can be integrated with the root filesystem using
mount points \parencite{freebsd_dir_struct}. The subdirectories \texttt{/cdrom}, \texttt{/mnt} \texttt{/tmp}, \texttt{/usr}, and \texttt{/var}
serve as common mount points of the filesystem \parencite{freebsd_dir_struct}. The children filesystems and the mount points they are
mounted on are recorded in the \texttt{/etc/fstab} directory \parencite{freebsd_dir_struct}. The \texttt{namei()} function handles
pathname resolutions \parencite{freebsd_namei}. \texttt{namei()} usually returns a vnode, either locked or unlocked, of the resolved
path \parencite{freebsd_namei}. The \texttt{namei()} function is often initialized with \texttt{NDINIT()} \parencite{freebsd_namei}.
For instance, we can specify the operation \texttt{namei()} will perform by passing either \texttt{LOOKUP},
\texttt{CREATE}, \texttt{DELETE}, and \texttt{RENAME} to the \texttt{op} argument \parencite{freebsd_namei}.
The resolved directory and its vnode being initialized to \texttt{namei()} is pointed by the \texttt{ndp} attribute \parencite{freebsd_namei}.
This increments the reference count of the vnode \parencite{freebsd_namei}. Decrementing the reference count would need to be done
by either \texttt{vrele()} or \texttt{vput()} depending on whether the \texttt{LOCKLEAF} flag was set or not \parencite{freebsd_namei}.
FreeBSD performs file lookups with the \texttt{VOP\_LOOKUP()} function \parencite{freebsd_voplookup}. It specifically searches for a
component in a specified directory path \parencite{freebsd_voplookup}. From a component, it returns a pointer to a locked vnode
\parencite{freebsd_voplookup}. This can be one of the sources of deadlocks, assuming that the filesystem does not follow a strict
tree hierarchy, since this process is rather complicated and is very centralized \parencite{freebsd_voplookup}.

\begin{table}[!htb]
	\centering
	\begin{tabular}{|c|p{0.6\textwidth}|}
		\hline
		\textbf{Subdirectory} & \textbf{Description}                                                                                                   \\
		\hline\hline
		\texttt{/bin}         & Contains utilities for both single- and multi-user environments                                                        \\ \hline
		\texttt{/boot}        & Contains programs and configurations necessary for booting up the system                                               \\ \hline
		\texttt{/dev}         & Contains device nodes, which are entry points to devices that are connected to the system \parencite{freebsd_intro}    \\ \hline
		\texttt{/etc}         & Contains system configuration files and scripts                                                                        \\ \hline
		\texttt{/mnt}         & Serves as a temporary mount point                                                                                      \\ \hline
		\texttt{/proc}        & Mounts the \textit{process file system}, which displays the process table of the filesystem \parencite{freebsd_procfs} \\ \hline
		\texttt{/rescue}      & Contains programs needed for system recovery during emergencies                                                        \\ \hline
		\texttt{/root}        & Acts as a home directory for the \texttt{root} account                                                                 \\ \hline
		\texttt{/sbin}        & System programs for both single- and multi-user environments                                                           \\ \hline
		\texttt{/tmp}         & Contains temporary files that are not preserved on system reboot                                                       \\ \hline
		\texttt{/usr}         & Contains the user utilities and tools                                                                                  \\ \hline
		\texttt{/var}         & Contains multi-purpose logs and spool files                                                                            \\ \hline
	\end{tabular}
	\caption{Common Subdirectories of the Root Directory in UNIX Filesystems \parencite{freebsd_dir_struct}}
	\label{root-subdirs}
\end{table}

For NetBSD specifically, directory entry lookups are often accelerated with the use of \textbf{name lookup caches} \parencite{netbsd_namecache}.
Whenever a pathname resolution is performed, it first checks the name lookup cache for an entry. If an entry exists, it retrieves the
associated entry in the cache. If an entry does not exist, it creates a new entry in the name lookup cache \parencite{netbsd_namecache}.
Cache lookups are performed by calling the \texttt{cache\_lookup()} function \parencite{netbsd_namecache}. The result of the cache lookup
is stored in the \texttt{vvp} argument \parencite{netbsd_namecache}. \texttt{cache\_lookup()} returns 0 if it encounters a cache miss
\parencite{netbsd_namecache}. It stores the resolved vnode in \texttt{vpp} for a positive cache hit \parencite{netbsd_namecache}. On
a negative hit, \texttt{vpp} is set to \texttt{NULL} \parencite{netbsd_namecache}. The \textit{Least Recently Used (LRU)} policy
is adopted by the name lookup cache so that entries that are frequently retrieved continue to stay in the cache \parencite{netbsd_namecache}.
Operations on the name lookup cache table are handled by the \texttt{cache\_enter()}, \texttt{cache\_purge()}, and \texttt{cache\_purgevfs()}
functions \parencite{netbsd_namecache}. New entries are added to the cache table with \texttt{cache\_enter()} \parencite{netbsd_namecache}.
Existing entries are removed or flushed from the cache table with \texttt{cache\_purge()} \parencite{netbsd_namecache}.
\texttt{cache\_purgevfs()} can be used to flush an entire mounted filesystem in the table \parencite{netbsd_namecache}.

Given their shared UNIX and 4.4BSD ancestry, it is expected that both systems share the same directory structure rooted at \texttt{/}
directory. Both systems also adopt the \texttt{namei()} function as their primary pathname resolution routine and \texttt{VOP\_LOOKUP()}
to resolve individual pathname components. Name cache lengths in both systems are capped at 31 characters, which is set by the
\texttt{NCHNAMLEN} macro \parencites{netbsd_namecache}{freebsd_namecache}. When entries exceed the cap, they are not added in the
name cache. Lastly, both systems face the risk of deadlocks with their \texttt{VOP\_LOOKUP()} implementation.

\subsection{Disk Space Allocation and Free Space Management} \label{space-allocation-and-free-space} % done

As mentioned in Section \ref{file-structure-and-hierarchy}, FreeBSD's UFS2 and NetBSD's FFS basically have the same organization. To summarize,
the disk is partitioned into cylinder groups. Each cylinder in the cylinder group consists of a superblock,
which contains metadata about the filesystem, the inodes, which contains metadata about files, and the
data blocks, which contains the contents of the files described by the cylinder's inodes. The contents of
a file are directly allocated within the first 12 logical blocks in the 15 direct addresses listed by the
file's inode \parencite{oracle2010}. Whenever the entire contents of the file cannot fit within the 12
logical blocks allocated for that file, the 13th address in the direct addresses list points to an indirect
block \parencite{oracle2010}. The indirect block contains direct addresses to storage locations \parencite{oracle2010}.
Whenever the file contents exceed the capacity of the 12 logical blocks, the remaining contents will
be allocated to the addresses in the indirect blocks. The last two direct addresses points to a double and triple
indirect blocks respectively. A double indirect block contains addresses to indirect blocks, which then
points to direct addresses \parencite{oracle2010}. The same pattern applies to the triple indirect block,
which holds addresses pointing to double indirect blocks \parencite{oracle2010}. The 15th address is only
used on very large files, hence it is only allocated when needed. In both UFS2 and FFS, a \textit{cylinder group map}
is allocated between the superblock and the inode block \parencite{oracle2010}. The cylinder group map
keeps track of unused blocks in the cylinder group, which could be blocks not used as inodes or indirect address
blocks or freed storage blocks \parencite{oracle2010}. These unused blocks are tracked to prevent fragmentation
\parencite{oracle2010}. For FFS, one inode
is assigned to every $4 \times frag\_size$ bytes of data where $frag\_size$ is the fragment size of the filesystem,
which scales with the size of the filesystem \parencite{netbsd_newfs}. For instance, if
the filesystem size is less than 20M, one inode is created for every 2K of data \parencite{netbsd_newfs}.

With ZFS, the filesystem is hosted in a \textbf{ZFS storage pool} or \textbf{Zpool}, which is a set
of disks used by the filesystem \parencite{data_recovery_glossary}. A Zpool is composed of virtual devices
or \textit{vdevs} \parencite{openzfs_vdev}. A vdev represents a logical grouping of physical storage devices such
as hard disks and SSDs \parencite{openzfs_vdev}. According to \textcite{openzfs_vdev}
vdevs are arranged in a tree-structured hierarchy, with leaf vdevs,
top-level vdevs, and a root vdev. Leaf vdevs directly point to a physical storage device. Top-level vdevs are the
children of the root vdev, which can correspond to a single device or a logical group of leaf vdevs. The root vdev
aggregates all the top-level vdevs into a single logical group, which is the Zpool.
According to \textcite{openzfs_vdev}, vdevs can be one of the following types:

\begin{itemize}
	\item \textbf{Striped Disk}: consists of one or more physical storage devices that are striped together. Data can be
	      lost in this vdev if it fails.

	\item \textbf{Mirror}: stores the same data on multiple devices. Mirrors are often used for redundancy.

	\item \textbf{RAIDZ}: a variation of RAID-5 \parencite{openzfs_raidz}; provides fault tolerance by utilizing parity.
	      RAIDZ can range from RAIDZ1 to RAIDZ 3. RAIDZ3 requires more disks than RAIDZ1 and RAIDZ2 but provides fault tolerance
	      on more drives than the two RAIDZ levels.

	\item \textbf{dRAID}: short for \textit{Distributed RAID}; provides a more distributed parity, allowing for quicker
	      recoveries after a drive fault.
\end{itemize}

\noindent
vdevs can be of auxiliary types such as \textit{Spare}, \textit{L2ARC}, \textit{SLOG}, \textit{Special}, and \textit{Dedup},
which have their own special uses. For instance, a Spare vdev is often used to replace a failed drive from another vdev.
A \textit{ZFS dataset} is automatically created in the creation of a Zpool \parencite{pasek2025}. A ZFS dataset handles
the organization of data within the filesystem \parencite{45drives}. ZFS allows creation of multiple ZFS datasets with
their own configuration such as level of compression, mountpoints, and snapshots \parencite{pasek2025}. By default,
the size of a record is set to 128K \parencite{pasek2025}. Compression is also not enabled by default \parencite{pasek2025}.
However, as mentioned earlier, these can be customized depending on the workload. According to \textcite{pasek2025},
for a general-purpose file server, a record size of 128K is recommended. For media archiving and backup workloads,
a record size of 1M is recommended \parencite{pasek2025}. While for database and VM disk workloads, a record size
of 16K is suggested \parencite{pasek2025}.

ZFS makes use of space maps to manage free disk space. The vdev is first partitioned into hundreds of chunks of
manageable size called \textit{metaslabs} \parencite{david2012}. Each metaslab is associated with a space map. A
\textit{space map} logs the allocation and deallocation activity \parencite{david2012}. The filesystem reads from
the logs to simulate the allocation activity into an AVL tree representing the free space of a vdev \parencite{david2012}.
At the same time, the filesystem condenses the space map by cancelling out allocation-free pairs \parencite{david2012}.
This way, the contiguous regions of memory that are currently allocated is tracked \parencite{siebenmann2015}.
With a space map, the filesystem can perform random frees just as fast as sequential frees since the amount of disk
space to be freed can be determined by the location of the metaslab listed in the space map offsetted by an extent
\parencite{david2012}.

UNIX-based filesystems such as UFS and FFS adopt the same bitmap approach for keeping track of free space. The bitmaps
serve as a primary reference to UFS and FFS to mitigate fragmentation. While bitmaps have minimal overhead, they often
cannot prevent fragmentation for large files. This typically happens when allocating files that exceed the 48K threshold
\parencite{freebsd_filesystem_org}. FFS resolves this by redirecting the block allocation to a different cylinder group
when it exceeds 48K and a megabyte after if it still exceeds the capacity of the redirected cylinder group \parencite{freebsd_filesystem_org}.
This adds complexity in allocating new blocks and can potentially increase seek distances when allocating large files.
ZFS, however, make use of space maps for free space management, which adopts a fundamentally different approach for managing free memory.
As discussed earlier, the space map approach significantly improves the performance of random frees since each metaslab
or a region in memory has its own space map wherein transactions are only logged on that region. For UFS and FFS,
random frees typically requires scattered bitmap updates, which can be inefficient as opposed to tracking free space from
transaction logs in the space map. However, this approach of logging the transactions and replaying them using AVL tree
can add significant overhead, which is nonexistent in the bitmap approach. Another advantage of ZFS is that it handles
well for large storage workloads since it is able to add more devices to the pool. This allows space allocation to be
redistributed across multiple vdevs. However, UFS and FFS are fixed into their partitions \parencite{schaumann},
hence expanding the partitions require the system to perform a full backup-repartition-recovery cycle. Free space management
is highly crucial in UFS as accessing files becomes increasingly inefficient as more files are added to the filesystem.
A good rule of thumb is to keep around 1-10\% of minimum free space, which depends on the partition size \parencite{oracle_min_freespace}.

\subsection{File Protection and Security} \label{file-protection-security} % done

Both FreeBSD and NetBSD extend \textit{owner/group/other} permission model for file access standard in UNIX systems with
POSIX's Access Control Lists. An \textbf{Access Control List (ACL)} extends the permission model by incorporating
a per-file or per-directory approach of access control \parencite{freebsd_security}. Each entry in the ACL consists of a user or group with
the associated permissions for a file or directory \parencite{freebsd_security}. FreeBSD offers the \texttt{getfacl(1)} and \texttt{setfacl(1)}
commands for managing the ACL \parencite{freebsd_security}. For FreeBSD, ZFS also introduces a more granular approach to ACLs,
which is the \textit{NFSv4 ACL} \parencite{oracle_acl2010}. The NFSv4 implementation of the ACL is more reminiscent to
NT's approach and offers a richer set of access permissions and inheritance semantics in declaring how these parmissions are
applied \parencite{oracle_acl2010}. For NetBSD, ACLs along with \textbf{extended attributes (EA)} are supported on later versions
of NetBSD, starting at NetBSD 10 \parencite{netbsd_ea_acl}. Furthermore, only the syntax of the ACL is known to the vnodes,
which are representations of files and directories on disk in FFS \parencite{netbsd_acl}. It is also possible for a file or
directory to be associated with zero or multiple ACLs \parencite{netbsd_acl}.

FreeBSD extends the read/write/execute permissions with \texttt{setuid}, \texttt{setgid}, and \texttt{sticky} permissions
\parencite{freebsd_permissions}. \texttt{setuid} replaces the \texttt{x} bit of the owner's permissions with \texttt{s},
which essentially allows the owner to use utilities that need higher privileges \parencite{freebsd_permissions}. This is often
enabled when changing user passwords with \tettt{passwd()} \parencite{freebsd_permissions}. \texttt{setgid} performs the same
operation as \texttt{setuid} on the group's permissions instead of the owner's \parencite{freebsd_permissions}. \texttt{sticky}
only allows file deletions by the file owner \parencite{freebsd_permissions}. This permission is usually set on a directory
\parencite{freebsd_permissions}. Additionally, \textit{file flags} extends the set of permissions only for files \parencite{freebsd_basics}.
They can be modified using the \texttt{chflag} command \parencites{freebsd_basics}{freebsd_chflags}. File flags include
the \textit{undeletable} flag which can be set by \texttt{sunlink} to prevent a file to be deleted even by the root user,
the \textit{user append-only} flag which can be set by \texttt{uappend} to only allow the user to perform file appends,
and the \textit{user immutable} flag which can be set by \texttt{uchg} to make the file immutable \parencite{freebsd_chflags}.
FreeBSD also offers \textit{secure levels} that can be adjusted with the \texttt{kern.securelevel} state variable
\parencite{freebsd_securelevel}. Table \ref{securelevel} lists the different secure levels that can be set by the user according to
\textcite{freebsd_securelevel}. By default, \texttt{securelevel} is set to \texttt{-1}.

\begin{table}[!htb]
	\centering
	\begin{tabular}{|c|c|p{0.5\textwidth}|}
		\hline
		\textbf{Secure Level} & \textbf{Mode}                      & \textbf{Description}                                                                                                                                                                                         \\
		\hline\hline
		\texttt{-1}           & \textit{Permanently insecure mode} & The system is in insecure mode.                                                                                                                                                                              \\ \hline
		\texttt{0}            & \textit{Insecure mode}             & Permissions are enforced when performing read/write operations to devices. Immutable and append-only flags can also be disabled.                                                                             \\ \hline
		\texttt{1}            & \textit{Secure mode}               & Immutable and append-only flags cannot be disabled. Specific mounted filesystems such as \texttt{/dev/mem}, \texttt{/dev/kmem}, and \texttt{/dev/io} cannot be accessed or only accessible with restrictions \\ \hline
		\texttt{2}            & \textit{Highly secure mode}        & Same restrictions in secure mode are enforced in this level. Additionally, mounted disks cannot be written. Kernel time changes are also restricted to within a second.                                      \\ \hline
		\texttt{3}            & \textit{Network secure mode}       & Same restrictions in highly secure mode are enforced in this level. Additonally, IP packet filter rules and \texttt{dummynet()} or \texttt{pf()} configuration cannot be modified.                           \\ \hline
	\end{tabular}
	\caption{Secure Levels in FreeBSD \parencite{freebsd_securelevel}}
	\label{securelevel}
\end{table}

In NetBSD, the permissions of a file can be checked with \texttt{VOP\_ACCESS()} \parencites{netbsd_vnodeops}{netbsd_internals}.
\texttt{VOP\_ACCESS()} can check if a file has the read, write, or execute permissions by passing in the \texttt{VREAD}, \texttt{VWRITE},
or \texttt{VEXEC} constants to the \texttt{mode} argument respectively \parencite{netbsd_vnodeops}. NetBSD also adopts the
\textbf{Role-Based Access Control (RBAC)} introduced by \textcite{crooks2009}. The RBAC model splits the privileges of a super user into
a total of 57 separate roles \parencite{crooks2009}. NetBSD also supports exploit mitigation mechanisms, which are the following according
to \textcite{netbsd_security}:

\begin{itemize}
	\item \textbf{PaX ASLR (Address Space Layout Randomization)}: obfuscate the address space layout, specifically where libraries
	      and specific functions are stored in memory.

	\item \textbf{PaX MPROTECT}: enforces memory protection by preventing malicious programs from marking writable memory blocks
	      as executable.

	\item \textbf{PaX SegvGuard}: tracks segmentation faults and provide rate-limiting.

	\item \textbf{\texttt{gcc()} Stack-Smashing Protection (SSP)}: aborts programs with corrupted variables and buffers which can be
	      used to alter the control flow of the program. These set of compiler extensions come with \texttt{gcc()}.

	\item \textbf{\texttt{FORTIFY\_SOURCE}}: detects buffer overflows with little to no bounds checking implemented to prevent
	      further damage.

	\item \textbf{NULL pointer dereferences protections}: restricts mappings on the zero address to prevent code injection
	      attacks.
\end{itemize}

\noindent
FreeBSD also shares similar exploit mitigation mechanisms with ASLR but only on 64-bit architectures \parencite{larabel2022}.
However, this is only adopted in the user address space as FreeBSD currently does not implement ASLR on the kernel space with
KASLR on FreeBSD 14.x releases \parencite{belousov_maste}.

Both systems adopt the same permission model standard to UNIX systems and extend the three standard permissions with POSIX ACLs.
Both systems also integrate EAs for their UNIX filesystems. For FreeBSD specifically, the EA layer is pluggable, which can provide
future flexibility at a cost of performance overhead. NetBSD's EAs are more directly integrated into the vnodes themselves,
however these mechanisms are only adopted in more recent releases of NetBSD, specifically the 10.0 release as mentioned earlier.
FreeBSD and NetBSD primarily differ in how permissions are being handled and extended. FreeBSD supports an extended set of permissions
in the form of file flags. Systemwide access permissions can also be configured with the \texttt{kern.securelevel} state variable,
which can lack granularity. In NetBSD, superuser permissions are split into permission roles through RBAC. While this provides an increased
precision in how permissions are granted, this designation of permissions into roles can be complex to manage and might require overhead.
In terms of exploit mitigation, NetBSD offers a wider variety of mitigation techniques with its PaX features and other standard
mitigation techniques such as enforcing restrictions to prevent intentional buffer overflows and NULL pointer dereferences.
However, these techniques need to be tuned for each program using \texttt{paxctl} to be compatible with code generators such as
debuggers and runtime environments such as the JVM \parencite{netbsd_paxctl}. This is contrast to FreeBSD's mitigation techniques
which are simpler to enforce but is less effective in actually mitigating exploits.

\subsection{Reliability and Crash Recovery} \label{reliability-and-crash-recovery} % done

In FreeBSD's UFS2 and NetBSD's FFS in the pre-6.0 releases, a \textbf{soft updated dependency tracking system} is employed for crash
recovery, serving as an alternative
to the journaled-filesystem approach \parencites{ganger1994}{mckusick1996}{netbsd_6.0}.
Soft updates ensure the integrity of file system metadata in an event of a crash by eliminating inconsistencies in the filesystem
on-disk \parencite{mckusick_ganger1999}. This is done by preventing metadata writes that renders the filesystem inconsistent and
enforcing rollbacks on metadata blocks that rely on potentially non-flushed blocks on write \parencite{mkcusick_ganger1999}.
With this approach, the operating system only needs to perform a background walk on the filesystem to free any orphaned space
\parencite{mckusick_ganger1999}. \textcite{mckusick2012} proposed an improvement to the
soft-update tracking system by implementing a journaling mechanism that keeps track of the inconsistencies produced by the
soft-update tracking system \parencite{mckusick2012}. This eliminates the need for a recovery operation after a crash, which
can be expensive \parencite{mckusick2002}.
In newer releases of NetBSD, FFS adopts \textbf{Write Ahead Physical Block Logging (WAPBL)} as its primary metadata journaling mechanism
for crash recovery \parencites{biancuzzi2010}{mcneill2008}. This eliminates the need to run \texttt{fsck} to check for any
inconsistencies on system boot in an event of a crash \parencite{netbsd_wapbl}. Instead, the operating system only needs to
simulate the logged transactions and correct any encountered inconsistencies \parencite{netbsd_wapbl}.

FreeBSD's ZFS employs multiple strategies to ensure data integrity and aid in crash recovery. For instance, it implements
\textbf{copy on writes (COW)}, which creates \textit{snapshots} of the same file on every write \parencites{pasek2025}{freebsd_zfs}.
Snapshots only keep track of the difference between each version of the file, which makes them space-efficient \parencite{freebsd_zfs}.
Snapshots also provide a way of making backups of the filesystem \parencite{freebsd_zfs}. Considering the efficiency that
snapshots bring, it makes it faster to restore previous snapshots than whole backups \parencite{freebsd_zfs}. Data integrity is
performed using \textbf{checksums}. Each data is associated with its checksum, which is generated when the data is stored in the disk
\parencite{freebsd_zfs}. The filesystem then recalculates the checksum on every read and performs error corrections using
redundant data if the recalculated checksum does not match with the checksum associated with the data \parencite{freebsd_zfs}.
Redundant data is produced with \textbf{RAIDZ}, which utilizes parity to create redundant copies of the same data across
different disks \parencite{freebsd_zfs}. As mentioned in Section \ref{file-structure-and-hierarchy}, RAIDZ has three different levels: RAIDZ1,
RAIDZ2, and RAIDZ3. A higher level of RAIDZ corresponds to a higher level of parity at the expense of having more disk requirements
\parencites{freebsd_zfs}{openzfs_vdev}. ZFS also implements transaction logging using the \textbf{ZFS Intent Log}, which
logs write transactions and stores them to the main pool devices \parencite{khurshid2025}. This allows the filesystem to replay the
logged transactions during system recovery \parencite{khurshid2025}.

The use of a soft updates dependency tracking for FreeBSD's UFS avoids overhead incurred from writing a journal, however, lost blocks
need to be reclaimed with \texttt{fsck}, which is known to be resource-intensive and time-consuming. This is unlike ZFS's ZIL in FreeBSD
and FFS's WAPBL in NetBSD which offer faster recovery at the expense of journal write overhead. Furthermore, ZIL's transaction
semantics is more sophisticated but it introduces a single point-of-failure that renders an entire pool to be corrupted when the
transaction logs themselves are corrupted. This problem is nonexistent with WAPBL due to its simpler design. However, it offers a lesser
range of recovery options. Data integrity is adopted in ZFS through checksums and RAIDZ redundancy, which is not natively employed
in FFS.

\subsection{Disk Scheduling and I/O Management} \label{disk-scheduling-io} % done

Earlier releases of FreeBSD used primitive disk schedulers such as elevator and C-LOOK schedulers \parencite{rizzo2009}.
\textcite{rizzo2009} later proposed utilizing the \textbf{GEOM architecture} in disk scheduling, which posed several advantages
such as ease of implementation of disk schedulers and reusability of existing mechanisms for locking and configuration.
GEOM uses a graph of GEOM modules to perform arbitrary manipulations of disk I/O requests or \textit{bio} \parencite{rizzo2009}. Disk I/O
requests enters from the provider nodes, which are connected to the filesystem, and exit to the consumer nodes, which
goes straight to the device driver \parencite{rizzo2009}. By default, GEOM implements a simple FIFO disk I/O scheduler
\parencite{losh2015}. The GEOM layer is connected to a \textbf{Common Access Method (CAM)} subsystem, which handles access of
media devices \parencite{losh2015}. Each disk driver below the GEOM layer implements their own scheduling policies
\parencite{losh2015}. CAM can support SCSI- and ATA-based disks with its two peripheral drivers for each of the disks
\parencite{losh2015}. Both peripheral drivers schedule disk I/O requests by scheduling as many \texttt{BIO\_READ}, \texttt{BIO\_WRITE},
or \texttt{BIO\_FLUSH} commands as it can \parencite{losh2015}. However, if it encounters a \texttt{BIO\_DELETE} command
and no other command is currently active, it combines as much \texttt{BIO\_DELETE} commands as it can and schedules them.
While the scheduling policies adopted on the GEOM layer and the CAM subsystem is sufficient for most of its users,
\textcite{losh2015} and the people working in Netflix discovered weaknesses to these algorithms, which include uneven I/O latency
if the system has long-read latency tails, no enforcement of priority, long latencies as a result of nonexistent deadlines,
and no way of balancing writing performance and reading latency.

In NetBSD, disk I/O requests are managed by a \textbf{buffer queue subsystem or \texttt{bufq}}, which manages device buffer
queues \parencite{netbsd_bufq}. The buffer queue subsystem provides several functions for manipulating the device buffer
queues, which includes \texttt{bufq\_put()} for adding an I/O request to the buffer queue and \texttt{bufq\_get} for retrieving
an I/O request from the buffer queue \parencites{netbsd_bufq}{chung2002}. NetBSD supports four buffer queue scheduling
strategies: \texttt{BUFQ\_FCFS}, \texttt{BUFQ\_DISKSORT}, \texttt{BUFQ\_PRIOCSCAN}, and \texttt{BUFQ\_READPRIO}
\parencites{netbsd_bufq}{netbsd_sysctl}{netbsd_options}. \texttt{BUFQ\_FCFS} is a simple implementation of the FCFS disk
scheduling policy. \texttt{BUFQ\_DISKSORT} is a one-way variant CSCAN disk scheduling policy by scheduling ascending
disk I/O requests. \texttt{BUFQ\_CSCAN} implements a CSCAN scheduling policy \parencite{netbsd_options}. \texttt{BUFQ\_READPRIO}
prioritizes disk read requests over disk write requests, which can be an ideal options for systems that frequently receives
disk write requests in one disk locality \parencite{netbsd_options}. By default, the \texttt{BUFQ\_FCFS} and \texttt{BUFQ\_DISKSORT}
are implemented in the kernel \parencite{netbsd_bufq}. When an I/O operation from a buffer queue is completed, the \texttt{biodone()}
function is called \parencite{chung2002}. NetBSD also implements a \texttt{biowait()} function that waits for an I/O operation
in a buffer queue to be completed \parencite{chung2002}. This is usually called inside a \texttt{bread()} or \texttt{bwrite()}
functions which allocate a buffer and release a modified buffer respectively \parencite{chung2002}.

GEOM's modular design in FreeBSD offers immense flexibility for implementing complex disk schedulers on the GEOM layer. This
flexibility does not seem to be observed in NetBSD's buffer queue approach. However, buffer queues in NetBSD connect directly
to the vnode layer, which can result in minimal overhead. Disk scheduling in FreeBSD is managed by the CAM subsystem, which
prioritizes \texttt{BIO\_READ}, \texttt{BIO\_WRITE}, and \texttt{BIO\_FLUSH} commands and combines multiple \texttt{BIO\_DELETE}
and scheduled as one. Scheduling mechanisms are decided by the lower-layer drivers, which can introduce unpredictability
in mixed workloads. However, NetBSD offers more explicit scheduling options with \texttt{BUFQ\_FCFS}, \texttt{BUFQ\_DISKSORT},
\texttt{BUFQ\_PRIOCSCAN}, and \texttt{BUFQ\_READPRIO}.

\subsection{Performance Considerations} \label{performance-considerations} % done

FreeBSD's ZFS supports three caching strategies: ARC, L2ARC, and SLOG \parencite{khurshid2025}. The
\textbf{ARC} or the \textbf{Adaptive Replacement Cache} is an adaptive read caching mechanism that resides in the system's main
memory \parencite{khurshid2025}. It aims to balance between recency of cache entries and frequency of access to the entries
\parencite{khurshid2025}. It competes with the kernel and user applications for memory \parencite{khurshid2025}. ARC balances
between the four caching lists according to \textcite{khurshid2025}:

\begin{itemize}
	\item \textbf{Most Recently Used (MRU)}: keeps a list of data that is not frequently accessed

	\item \textbf{Most Frequently Used (MFU)}: keeps a list of data that is frequently accessed

	\item \textbf{Ghost MRU}: contains metadata on the recently removed entries in the MRU cache list

	\item \textbf{Ghost MFU}: contains metadata on the recently removed entries in the MFU cache list
\end{itemize}

Both Ghost MRU and Ghost MFU do not cache actual data, but informs ARC on the removed cache entries, which allows the algorithm
to adjust the memory allocation for new entries \parencite{khurshid2025}. This self-tuning mechanism increases the frequency of
future cache hits, which makes the algorithm highly efficient for a variety of workloads \parencite{khurshid2025}. A similar
caching mechanism is also applied when accessing fast secondary disk devices with the \textbf{L2ARC} cache \parencite{khurshid2025}.
The L2ARC cache typically stores data that has been evicted or will eventually be evicted by the main cache \parencite{khurshid2025}.
Instead of actually flushing the entry, ARC keeps a pointer to where the evicted memory block is stored in the L2ARC device
\parencite{khurshid2025}. This makes accesses to recently evicted entries more efficient. It is possible for multiple L2ARC devices
to be added to the system, as long as there is sufficient main memory to store the block pointers without compromising the speed
of the system due to too many block pointers in RAM \parencite{khurshid2025}. For write operations, the \textbf{Separate Log Device (SLOG)}
is often used \parencite{khurshid2025}. SLOG increases the latency for synchronous write requests considering that they need to
be acknowledge right away \parencite{khurshid2025}. However, it does not improve the disk request throughput. \parencite{khurshid2025}.
Hence, SLOG can achieve an improved performance for low-latency disk devices \parencite{khurshid2025}. ZFS also allows for
performance-tuning. For instance, the memory usage for the ARC cache can be tuned with the \texttt{arc\_max} and \texttt{arc\_min}
tunables \parencite{khurshid2025}. Frequency of cache writes to the L2ARC device can be adjusted with the \texttt{l2arc\_feed\_min\_ms}
and \texttt{l2arc\_feed\_secs} tunables \parencite{khurshid2025}. The \texttt{l2arc\_write\_max} tunable can also be modified to
adjust the maximum cache write rate of the L2ARC device, which can avoid wearing out the device from writing data that will
not be used to the cache \parencite{khurshid2025}. The filesystem itself can also be tuned to increase either latency or throughput
with the \texttt{logbias} property \parencite{lucas2015}. By default, \texttt{logbias} is set to \textit{latency} to ensure that the data
is quickly synced to the disk \parencite{lucas2015}. However, it can be set to \textit{throughput}, which can be beneficial when
accessing large files \parencite{lucas2015}. The \texttt{vfs.zfs.prefetch.disable} can also be set to either enable or disable
prefetch by setting \texttt{vfs.zfs.prefetch.disable} to \texttt{1} or \texttt{0} respectively \parencite{freebsd_zfs}. Prefetching
allows bigger ARC blocks to be read than requested, which can lead to decreased performance if the workload contains a lot of
random reads to the cache \parencite{freebsd_zfs}.

NetBSD's FFS unifies its filesystem and virtual memory caches into a single subsystem called the \textbf{Unified Buffer Cache (UBC)}
\parencite{silvers2000}. The UBC addresses the inefficiency and poor flexibility brought by the two separate caching mechanisms in
the VM and I/O subsystems \parencite{silvers2000}. UBC stores accessed data from \texttt{read()/write()} and \texttt{mmap()} calls
to a page cache \parencite{silvers2000}. Read operations retrieve data from the page cache instead of the buffer cache \parencite{silvers2000}.
UBC also maps pages in the page cache to its corresponding address in memory \parencite{silvers2000}, as previously mentioned in
Section \ref{file-access-methods}. Page entries can be added and retrieved with the \texttt{VOP\_GETPAGES()} and \texttt{VOP\_PUTPAGES()}
functions respectively \parencite{silvers2000}. Additionally, \texttt{VOP\_GETPAGES()} needs to allocate data that is not yet stored
in the cache and \texttt{VOP\_PUTPAGES()} needs to write the dirty pages back to the disk; both requires the initiation of an I/O
operation \parencite{silvers2000}.

The caching approaches of ZFS and FFS vary in their complexity. ZFS's ARC is a deeply layered system with MRU/MFU cache lists, ghost
tracking and L2ARC for retrieving recently evicted entries, which can be fine-tuned with tunables such as \texttt{arc\_max},
\texttt{l3arc\_feed\_secs}, and \texttt{logbias}. ARC caching strategy is ideally used in high-storage workloads as it boosts
storage access performance. However, this deeply layered system introduces high complexity, which can actually hurt its performance
instead of boosting it. For instance, L2ARC adds RAM overhead for tracking block pointers. Enabling
prefetching can degrade performance on workloads that perform a high rate of random-read requests. Careless tuning of the tunables
can also negatively impact the performance and stability of the system \parencite{khurshid2025}. Unlike ZFS, FFS employs a
simpler and unified caching mechanism with UBCs, which often makes it an ideal option for general-purpose workloads.
This simplicity comes at a cost of lesser flexibility and difficulty in tuning. Furthermore, early benchmarks on UBC tested by
\textcite{silvers2000} show that the write performance of UBC still needs improvement.

\subsection{Summary} \label{file-and-disk-summary}

FreeBSD and NetBSD largely diverge in the design philosophy of the major filesystems they adopt and implement. This divergence
also defines the core strengths and weaknesses of each of their implementations. FreeBSD's ZFS unifies the filesystem and volume
manager into a single system built on a layered dnode/object set/DSL hierarchy, storage pools composed of vdevs, free space
management with space maps, and a layered caching subsystem with ARC, L2ARC, and SLOG. This design contributes to the reliability
and scalability that FreeBSD offers. In addition to this core design, ZFS also comes with additional features for data integrity
such as copy-on-write snapshots, checksums with error correction, and RAIDz parity. The ability to add more vdevs in a storage
pool allows for storage allocation to be redistributed across disk devices without needing an entire backup-repartition-recovery
cycle. This design comes at a cost of complexity and resource overhead. Using space maps for free space management incurs overhead
in the form of transaction logging and AVL-replays. ARC can often compete with the kernel and user applications for memory space.
L2ARC adds overhead on the RAM for storing the block pointers. Careless tuning can even hurt the performance of the system. In
contrast, NetBSD relies on a simpler implementation with FFS, which retains the classic cylinder-group organization consisting of
superblocks, inodes, bitmaps for tracking free space, and data blocks. FFS also adopts a unified caching mechanism with UBC.
These design choices contributes to NetBSD's high cross-compatibility on a wide variety of platforms and workloads. However, it
lacks the additional features that ZFS provides such as checksumming and RAIDZ redundancy. Partitions are also fixed in size,
and expanding them is an expensive process. Its free space management mechanism with bitmaps cannot prevent fragmentation when
allocating file sizes over 48K.

The tradeoffs between each design philosophy is evident across the four comparison dimensions. In terms of performance, ZFS's
adaptive and self-tuning ARC can boost access for high-storage workloads, but its complexity can add significant overhead.
Careless tuning can also hurt its performance. NetBSD's UBC is simpler and well-suited for general-purpose workloads, though
early benchmarks suggest that its write performance needs significant improvement. It also offers less flexibility in tuning.
In terms of reliability, both initially adopted soft updates for metadata consistency, but on more newer releases of FreeBSD
adopting ZFS offers more sophisticated recovery mechanisms with copy-on-write, checksumming, RAIDZ, and the ZIL. NetBSD
incorporates its own journaling mechanism with WAPBL, which is much simpler and has no single point of failure but has less recovery options.
This is unlike ZIL which can corrupt an entire storage pool if the logs are corrupted. In terms of scalability, ZFS's pool-and-vdev
model allows for greater scalability by allowing expansion of allocated memory across devices. UFS and FFS does not currently
have such mechanisms. In terms of security, the two systems share the standard UNIX permissions model which are extended
with POSIX ACLs and EAs. However, the two operating systems differ in how they approach extended permissions. On one hand,
FreeBSD allows for more refined permissions with file flags and a more coarse-grained and systemwide permissions control
with secure levels. On the other hand, NetBSD splits the superuser privileges into 57 RBAC roles. Furthermore, NetBSD offers
an extensive set of mitigation techniques with PaX features at a cost of per-program tuning complexity.

These tradeoffs show which system is ideal under specific conditions and workloads. For workloads that perform frequent
file writes such as a logging system, NetBSD's FFS ensures data corruption by
committing metadata logs and handles predictable I/O with minimal overhead with its buffer queue scheduling strategies.
FreeBSD's ZFS would absorb writes through copy-on-write and the ZIL, which can be accelerated by a SLOG device for low-latency
synchronous writes. For a system handling large multimedia files, FreeBSD's ZFS would be a more preferable option with its
record sizes reaching 1M capacities for media and backup workloads. Furthermore, its \texttt{logbias} property could be
configured to prioritize throughput over latency. It also allows for scalability with its pooled vdevs which can accommodate
large datasets. In contrast, NetBSD's FFS might degrade in performance because of its increased fragmentation when file
sizes exceed 48K, which can add allocation complexity and increasing seek distances. Overall, FreeBSD offers high data integrity,
scalability, and performance for intense workloads at a cost of high complexity and overhead. This makes FreeBSD an ideal
option for specific large-scale and high-intensity workloads that could greatly benefit from ZFS's complexity such as servers,
backup systems, and cloud storage. While NetBSD offers simplicity, portability, and predictability at a cost of a limited set of
additional features. This makes NetBSD an ideal option for general-purpose and resource-constrained workloads such as desktops,
laptops, and IoT systems.
